\section{AgentQ：从“失败截图”到可学习的恢复轨迹}

这一段是整堂课最容易被 demo 节奏掩盖的部分。讲者先展示常见失败，再展示预约、日历和航班任务。真正的教学点不是“agent 能订餐厅”，而是系统怎样搜索替代路径、修正错误、同时维护多个决策，并把成功与失败轨迹回收为训练数据。

\subsection{先承认 Agent 并不完美}

失败 collage 包含表单校验、支付、不可用时间和错误页面。它提醒读者：真实环境的反馈往往不是一个干净的 reward，而是错误消息、空结果、半完成页面和隐藏约束。AgentQ 的研究动机正是让 agent 在这些状态下继续搜索，而不是把一次失败当作任务结束。

这一小节承接上一章的评测分布：REAL 告诉我们失败集中在哪些站点，AgentQ 则追问失败之后怎样行动。阅读开场图时应先把错误按“可恢复”和“不可恢复”分开：表单字段缺失通常可修改，支付被拒可能需要用户介入，错误身份或错误收件人则可能必须终止。只有先定义恢复边界，搜索算法才不会把“坚持尝试”误当成正确行为。

\lecturefigure{slide-19-agentq-failure-collage.jpg}{AgentQ 开场先展示真实网页失败类型}{00:17:59}

\lecturefigure{slide-20-agentq-title.jpg}{AgentQ：面向自主 Agent 的搜索与学习框架}{00:18:08}

\begin{importantbox}{失败必须被结构化}
一次失败至少分为 observation error、planning error、tool/action error、environment constraint 与 policy/safety rejection。只记录“任务失败”无法决定应该改 prompt、工具、数据、奖励还是权限。
\end{importantbox}

\subsection{餐厅任务：搜索、回退与完成}

预约 demo 的关键转折是目标时间不可用。Agent 先查 OpenTable，再用 Google 获取补充信息，最后返回原站点完成预约。下面四个状态不是为了展示动画，而是说明一个恢复过程：动作失败、跨工具搜索、重新定位、完成并确认。

\lecturefigure{slide-21-reservation-prompt.jpg}{用户目标：指定餐厅、时间与人数}{00:18:11}

\lecturefigure{slide-22-agent-web-actions.jpg}{Agent 在网页中执行预约动作}{00:18:15}

\lecturefigure{slide-23-google-fallback.jpg}{原站点失败后转向 Google 获取可用性信息}{00:18:19}

\lecturefigure{slide-24-adaptable-site-formats.jpg}{跨站点返回并适配不同页面结构}{00:18:22}

\lecturefigure{slide-25-reservation-confirmed.jpg}{任务完成：预约确认成为可验证后置条件}{00:18:25}

\begin{knowledgebox}{读图：恢复链应该怎样验收}
先确认 agent 是否识别“10 PM 不可用”这一环境约束；再看跨站搜索是否引入了同名餐厅、日期和时区混淆；最后检查确认页是否与原始意图一致。只有最终页面、结构化订单和用户目标三者对齐，任务才算成功。
\end{knowledgebox}

\subsection{Self-correction 与并行决策}

下一组图把能力拆成三个层级：能够自我修正、能够同时处理多个决策、能够识别不合理日历时间并重排。这里的“self-correction”不应理解成模型说一句“我错了”，而是状态机检测后置条件失败、生成替代动作并保留原任务约束。

\lecturefigure{slide-26-self-correcting-behavior.jpg}{Agent 行为树：失败分支后重新选择路径}{00:18:28}

\lecturefigure{slide-27-multiple-decisions.jpg}{同一任务中维护多个并行决策}{00:18:31}

\lecturefigure{slide-28-calendar-prompt.jpg}{日历任务：明天 3 点安排一小时会议}{00:18:32}

\lecturefigure{slide-29-calendar-rescheduled.jpg}{识别 3 AM 不合理并重排到 3 PM}{00:18:39}

\begin{warningbox}{“自动修正”最容易越权}
如果用户写“3”，系统把它从 3 AM 改成 3 PM 是合理推断还是擅自改变？低风险日历任务可以通过确认解决；高风险任务必须把修正方案呈现给用户，而不是静默执行。
\end{warningbox}

\subsection{Skills routing 与动作上下文}

“right skills at the right time” 强调工具路由。航班任务需要解析城市、日期、单程、价格偏好和座位偏好，再选择搜索工具。路由器不只是分类器；它还要决定哪些约束进入工具参数，哪些保留给后续筛选，以及失败后是否换工具。

前面的餐厅与日历案例主要展示恢复，本节转向\textbf{能力选择}。一个 agent 即使拥有十个工具，也不代表应该把所有工具同时暴露给模型；工具越多，名称相似、参数冲突和权限误用越容易发生。读图时应关注约束在链路中的保真度：用户说“preferably one-way”和“window seat”，搜索阶段可以暂时放宽价格，但不能在最终推荐时忘掉座位偏好，更不能因为某个工具返回快就静默改变出发日期。

\lecturefigure{slide-30-right-skills-right-time.jpg}{正确时间调用正确 skill 是 Agent 编排问题}{00:18:41}

\lecturefigure{slide-31-flight-prompt.jpg}{航班请求包含日期、单程、价格与座位偏好}{00:18:42}

\lecturefigure{slide-32-flight-window-seat.jpg}{搜索结果中继续筛选经济舱靠窗座位}{00:18:48}

\lecturefigure{slide-33-booking-success-improvement.jpg}{课堂给出的 booking task success rate 提升}{00:18:54}

这一串 demo 也暴露了最终成功率之外的隐性成本。跨站搜索增加页面加载与 token，日历纠错需要解释时间歧义，航班筛选同时维护多个软约束；一个看似简单的“替我完成”请求，实际上包含状态跟踪、工具切换、偏好保持和结果验证。若系统只优化成功标签，可能学会更激进地点击或忽略软约束。训练数据必须把用户意图的完整保真度作为结果的一部分。

\teachervoice{00:17:58--00:20:47，讲者明确说 AgentQ 不是完美系统；这些 demo 用来展示失败、回退、自修正、多决策和 skill routing。讲义提醒：课堂结果属于特定任务与模型设置，不应推广为任意网站的可靠性。}

\begin{importantbox}{一个最小可复现 Agent trace}
\begin{lstlisting}[caption={Agent 轨迹应保存的最小字段}]
trace = {
  "goal": user_goal,
  "initial_state": state_snapshot,
  "steps": [
    {"observation": obs, "thought_summary": plan,
     "action": action, "tool_result": result,
     "postcondition": check, "cost": cost}
  ],
  "outcome": outcome,
  "user_confirmations": confirmations,
  "rollback": rollback_record
}
\end{lstlisting}
\end{importantbox}

\subsection{本章小结}

AgentQ 的 demo 不是“自动订票”宣传，而是一组失败恢复案例：环境给出约束，agent 搜索替代路径，状态机验证后置条件，路由器选择 skill，并把轨迹保存为学习材料。下一章把这些现象写成 MCTS、process feedback 和 preference learning。

